%% ============================================================
%% Academic Journal Paper Template
%% Compatible with: IEEE, ACM, Elsevier, Springer conventions
%% ============================================================

\documentclass[10pt, twocolumn]{article}

%% ---- Core Packages ----------------------------------------
\usepackage[T1]{fontenc}
\usepackage[utf8]{inputenc}
\usepackage{lmodern}
\usepackage{microtype}           % Microtypographic refinements

%% ---- Page Geometry ----------------------------------------
\usepackage[
  top=1in, bottom=1in,
  left=0.75in, right=0.75in,
  columnsep=0.25in
]{geometry}

%% ---- Mathematics ------------------------------------------
\usepackage{amsmath, amssymb, amsthm}
\usepackage{mathtools}
\usepackage{bm}                  % Bold math symbols

%% ---- Figures & Tables -------------------------------------
\usepackage{graphicx}
\usepackage{booktabs}            % Professional table rules
\usepackage{multirow}
\usepackage{array}
\usepackage{caption}
\usepackage{subcaption}
\usepackage{float}

%% ---- Algorithms -------------------------------------------
\usepackage[ruled, vlined, linesnumbered]{algorithm2e}

%% ---- Code Listings ----------------------------------------
\usepackage{listings}
\usepackage{xcolor}

\lstdefinestyle{codestyle}{
  backgroundcolor=\color{gray!8},
  basicstyle=\ttfamily\footnotesize,
  breakatwhitespace=false,
  breaklines=true,
  captionpos=b,
  commentstyle=\color{green!50!black},
  keywordstyle=\color{blue}\bfseries,
  stringstyle=\color{orange!70!black},
  numberstyle=\tiny\color{gray},
  numbers=left,
  numbersep=5pt,
  showstringspaces=false,
  frame=single,
  rulecolor=\color{gray!40},
  tabsize=2
}
\lstset{style=codestyle}

%% ---- Hyperlinks -------------------------------------------
\usepackage[
  colorlinks=true,
  linkcolor=blue!70!black,
  citecolor=green!50!black,
  urlcolor=blue!60!black
]{hyperref}
\usepackage{url}

%% ---- Bibliography -----------------------------------------
\usepackage[numbers, sort&compress]{natbib}

%% ---- Miscellaneous ----------------------------------------
\usepackage{lipsum}              % Lorem ipsum filler — remove in production
\usepackage{enumitem}
\usepackage{siunitx}             % SI units: \SI{3.14}{\mega\hertz}
\usepackage{cleveref}            % Smart cross-references: \cref{fig:foo}

%% ---- Theorem Environments ---------------------------------
\theoremstyle{plain}
\newtheorem{theorem}{Theorem}[section]
\newtheorem{lemma}[theorem]{Lemma}
\newtheorem{corollary}[theorem]{Corollary}
\newtheorem{proposition}[theorem]{Proposition}

\theoremstyle{definition}
\newtheorem{definition}[theorem]{Definition}
\newtheorem{example}[theorem]{Example}

\theoremstyle{remark}
\newtheorem{remark}[theorem]{Remark}

%% ---- Custom Commands --------------------------------------
\newcommand{\R}{\mathbb{R}}
\newcommand{\N}{\mathbb{N}}
\newcommand{\Z}{\mathbb{Z}}
\newcommand{\E}{\mathbb{E}}
\newcommand{\Prob}{\mathbb{P}}
\newcommand{\norm}[1]{\left\lVert #1 \right\rVert}
\newcommand{\abs}[1]{\left\lvert #1 \right\rvert}
\newcommand{\inner}[2]{\left\langle #1,\, #2 \right\rangle}
\newcommand{\ie}{\textit{i.e.}\xspace}
\newcommand{\eg}{\textit{e.g.}\xspace}
\newcommand{\etal}{\textit{et al.}\xspace}

%% ---- Caption Formatting -----------------------------------
\captionsetup{
  font=small,
  labelfont=bf,
  format=hang,
  justification=justified
}

%% ---- Header / Footer --------------------------------------
\usepackage{fancyhdr}
\pagestyle{fancy}
\fancyhf{}
\fancyhead[L]{\small\itshape Ne.org, Vol.~2, No.~1, March 2026}
\fancyhead[R]{\small\thepage}
\renewcommand{\headrulewidth}{0.4pt}

%% ============================================================
%%  FRONT MATTER
%% ============================================================

\title{%
  \vspace{-1.5em}%
  \large\textbf{The N-Token Concept.}%
}

\author{%
  \textbf{Amlal El Mahrouss}$^{1}$\thanks{Corresponding author. Email: amlalelmahrouss@icloud.com}
  \\amlalelmahrouss@icloud.com\\[0.4em]
}

\date{%
  \small 21 March, 2026
}

%% ============================================================
\begin{document}
%% ============================================================

\twocolumn[{%
  \maketitle
  \vspace{-0.5em}
  \rule{\linewidth}{0.4pt}
  %% ---- Abstract -------------------------------------------
  \begin{center}
  \begin{minipage}{0.92\linewidth}
    \small
    \textbf{Abstract.}\enspace
    This development paper presents an algorithm which expresses characters, 'A', 'B', 'Z' as single tokens. Each character has an index chosen by two parties, e.g: A=428, B=422.

    \medskip
    \textbf{Keywords:}\enspace
    Computer Science; Software Engineering;
  \end{minipage}
  \end{center}
  \vspace{0.8em}
  \rule{\linewidth}{0.4pt}
  \vspace{1em}
}]

%% ============================================================
\section{Introduction}
\label{sec:acode-introduction}

LLMs have for a long-time struggled with multi-width characters in their transformers architecture. This paper aim to address to current gaps present in such transformers by introducing an algorithm which divide multi-width characters into digestable chunks. 

\section{Theorem 1}
\label{sec:acode-theorem01}

We shall introduce the following algorithm of A-Indices.
Let the following algorithm for indices elements in an A-Base:
\begin{equation}
    \operatorname{AGroup}(x^{n-1}+x^{n-2}+x^{n})+\operatorname{AGroup}(x^{n+2}+x^{n+1})
\end{equation}
\begin{equation}
    \forall x \in \mathbb{R}, n > 2, x > 1, \forall n \in \mathbb{Z}
\end{equation}
Where $x$ is the A-Base index. An $\operatorname{AAlgorithm(0)}$ is an invalid index. Each element $x$ shall be divisible by a power of two. An $\operatorname{AAlgorithm}$ is an algorithm used to encode an index to `A-Code'.

\section{Theorem 2}
\label{sec:acode-theorem02}

Let a set of $\operatorname{AAlgorithm(x)}$ be elements of an `A-Code' alphabet. From a start index to an end index. An $\operatorname{AAlgorithm(x)}$ should be non-zero and continuous in the `A-Code' set. 

%% ============================================================
\section{Applications: The CJK characters}
\label{sec:acode-conclusion}

After laying the two theorems, we will use the following Python script:

\begin{lstlisting}
import pytoklib

if __name__ == '__main__':
    with open('datasets/jp.txt') as f: 
        s = f.read()
        res = pytoklib.ntoken.ntoken(s, pytoklib.sa1)
        for vi in res:
            print(vi[0])
            print(vi[1])

\end{lstlisting}

%% ============================================================
\section{Conclusion}
\label{sec:acode-conclusion}

Concluding the paper with this approach could be useful when dealing with large, multi-width characters.

\end{document}
